\section{Discussão}

O resultado primário indica que a ConvNeXt-Tiny discrimina células normais e
anormais da CRIC com AUC ROC de $0{,}9709\pm0{,}0059$ e AUC PR de
$0{,}9607\pm0{,}0121$ na validação cruzada agrupada. A estabilidade da AUC entre
folds sugere boa capacidade de ranqueamento, mas a sensibilidade variou de
0,8253 a 0,9434. Essa variação é central para a interpretação clínica: em uma
tarefa de triagem, a ordenação global pode ser alta e, ainda assim, a taxa de
células anormais recuperadas depender da composição das imagens avaliadas.
Essa é a principal entrega do estudo: transformar um resultado de
classificação em uma análise de uso, variância e erro, em vez de apresentar
apenas uma acurácia agregada.

A análise Bethesda oferece uma explicação plausível para parte dessa
variabilidade. Na partição retida, HSIL, categoria com 99,4\% de acerto, está
super-representada entre as anormais, enquanto LSIL está sub-representada e
ASC-US apresenta 81,7\% de acerto. Assim, partições com mais HSIL tendem a
parecer mais favoráveis ao classificador, ao passo que maior presença de LSIL e
ASC-US pode reduzir a sensibilidade. Como a composição por subtipo de cada fold
não foi modelada formalmente, essa ligação deve ser tratada como hipótese
apoiada pela análise estratificada, não como explicação causal demonstrada.

Na partição retida, a comparação com ResNet50 mostra que a ConvNeXt-Tiny foi
superior em AUC ROC e AUC PR, mas não dominou todas as métricas dependentes de
limiar. A ResNet50 alcançou maior especificidade e MCC na comparação local,
enquanto a ConvNeXt-Tiny apresentou maior sensibilidade e melhor ranqueamento.
Esse padrão reforça a separação entre duas perguntas: qual modelo ordena melhor
as células por risco e qual ponto operacional produz a melhor decisão binária
para um fluxo específico.

Em relação à literatura, a comparação deve usar a validação cruzada agrupada, e
não apenas a partição com resultado mais alto. Estudos em Herlev e SIPaKMeD
reportaram acurácias elevadas com ResNet50, CCanNet e ensembles
\cite{kaur2025comparison,mehedi2024lightweight,lima2025ensemble}, mas diferem
em preparação, classes, unidade de particionamento e métrica principal. O
resultado médio de AUC ROC 0,9709 na CRIC não estabelece superioridade nem
inferioridade diante dessas acurácias. Ele indica, de modo mais restrito, que
um protocolo agrupado por imagem em citologia convencional ainda permite
desempenho competitivo. O desempenho inferior reportado por Rodríguez et al. em
esfregaços convencionais reforça que essa preparação pode ser mais desafiadora
que a citologia em meio líquido \cite{rodriguez2024automated}.

A calibração também deve ser interpretada no escopo da partição retida. O
escalonamento por temperatura reduziu o ECE de 0,0247 para 0,0130, mas esse
ganho não foi avaliado por validação cruzada. Além disso, uma transformação
monotônica não altera AUC; a diferença de ranqueamento observada na análise
detalhada decorre do TTA. Em cenários de uso clínico, probabilidades calibradas
são relevantes para priorização e comunicação de incerteza, mas exigem
validação externa e amostras maiores.

\subsection{Recomendações para estudos comparáveis}

Uma contribuição adicional do estudo é explicitar escolhas de avaliação que
deveriam ser reportadas em trabalhos futuros com células cervicais recortadas.
Diretrizes recentes para modelos preditivos e IA em imagem médica recomendam
declarar o uso pretendido, a fonte dos dados, o nível de particionamento, as
métricas e sua incerteza, além das limitações de generalização
\cite{collins2024tripod,tejani2024claim}. No contexto da CRIC, esses itens são
particularmente importantes porque múltiplas células podem compartilhar a mesma
imagem, lâmina, coloração e artefatos de aquisição.

\begin{table}[ht]
\centering
\caption{Elementos mínimos recomendados para avaliação de triagem celular em
citologia cervical convencional.}
\label{tab:recomendacoes}
\scriptsize
\renewcommand{\arraystretch}{0.92}
\setlength{\tabcolsep}{3pt}
\begin{tabular}{p{3.1cm}p{4.25cm}p{4.35cm}}
\toprule
Elemento & Risco mitigado & Como foi tratado neste estudo \\
\midrule
Unidade de partição & Vazamento entre recortes visualmente correlatos &
Separação por imagem-fonte e validação cruzada agrupada \\
Métrica primária & Seleção oportunista da partição mais favorável & AUC ROC,
AUC PR e métricas de decisão reportadas como média $\pm$ desvio-padrão \\
Limiar operacional & Confundir ranqueamento com decisão clínica & Comparação
entre 0,50, Youden e regime de alta sensibilidade \\
Incerteza & Tratar células correlatas como independentes & Bootstrap agrupado
por imagem e intervalos de Wilson por subtipo \\
Análise por subtipo & Ocultar falhas em categorias limítrofes & Estratificação
por NILM, ASC-US, LSIL, ASC-H, HSIL e SCC \\
\bottomrule
\end{tabular}
\renewcommand{\arraystretch}{1.0}
\end{table}

Essa estrutura também ajuda a interpretar diferenças entre estudos. Uma
acurácia elevada em Herlev ou SIPaKMeD pode ser útil como referência de
arquitetura, mas não substitui um protocolo que declare o nível de separação
dos dados, o regime de decisão e a composição dos subtipos. Para a CRIC, a
ausência de identificadores de paciente torna inadequado apresentar o
particionamento por imagem como solução definitiva; ele deve ser descrito como
o melhor controle disponível na base, não como garantia clínica.

\subsection{Limitações e uso responsável}

A principal limitação de validade está no nível de independência dos dados. O
agrupamento por imagem reduz o compartilhamento direto de recortes entre treino
e teste, mas não garante separação por paciente, pois a CRIC não fornece
identificadores individuais e suas 400 imagens provêm de um número menor de
lâminas. Além disso, a ResNet50 foi treinada na mesma partição retida, mas não
foi repetida em todos os folds; portanto, a comparação entre arquiteturas deve
ser lida como análise complementar.

Também há limitações de escopo. A tarefa binária avalia recortes centrados em
núcleos anotados, não a leitura de uma lâmina completa. O preenchimento branco
em bordas pode funcionar como atalho visual, e recortes de regiões sobrepostas
podem conter núcleos vizinhos que influenciam a decisão atribuída à célula
central. Não houve validação externa em outra instituição, equipamento ou
população, o que restringe conclusões sobre generalização.

Por essas razões, o sistema deve ser entendido como protótipo de pesquisa para
priorização: ele recebe uma célula já localizada e devolve uma pontuação para
revisão. Um fluxo completo ainda necessita localizar células, lidar com campos
sem anotações, agregar evidências no nível da lâmina e incorporar supervisão
humana. A coleção CRIC é pública, anonimizada e distribuída para pesquisa
\cite{rezende2021cric}; o modelo não foi validado prospectivamente e não deve
ser usado para diagnóstico ou exclusão de casos.
